\begin{textsample}{POS Dim 5 – global\_north\_2024\_08 – Score 10.00 – global\_north\_2024\_08/111634104.txt}  \label{ex:f5_pos_380}
Aotearoa Healthcare Workers for Palestine ( AHW4P ) has written to all New Zealand medical institutions demanding that they break their silence over the Israeli \textbf{genocide} and attacks on healthcare in Gaza ( 1 ).
The open letter was signed by more than 300 healthcare workers from across the health sector.
It was endorsed by Ora Taiao : NZ Climate and Health Council, Maaori Doctors in Solidarity with Palestine, Te Ohu Rata O Aotearoa ( Maaori Medical Practitioners ), The Maaori Health Advisory Group of the Royal Australasian College of Surgeons and Te Kaahui Manukura O Kai Ora ( Maaori Dietitians Association ).
Hundreds of members of the public also signed the letter.
Whilst two Australasian medical colleges have since released statements, the vast majority of our medical institutions remain silent.
This is despite the fact that many have expressed solidarity over Russian attacks on healthcare in Ukraine."It is unconscionable that many of our medical colleges and professional associations are still silent after ten months, when an entire health system has been destroyed", says AHW4P spokesperson Grant Brookes. @ @ @ @ @ @ @ @ @ @ prestigious medical journals in the world, last month published a report of 186,000 deaths as a result of Israel ' s invasion ( 2 ).
The ICJ \textbf{rulings} are clear.
Israel is an apartheid state and is plausibly \textbf{committing} a \textbf{genocide} in Gaza.
Silence is not an option."Last week, 45 US healthcare workers who volunteered in Gaza wrote to the US government, demanding a ceasefire and arms \textbf{embargo} ( 3 ).
Advertisement - \textbf{scroll} to continue reading Are you getting our free newsletter?"The eyewitness testimony is devastating", Says Brookes."Some of these healthcare workers have vast experience in the field of humanitarian medicine and they are reporting that this is the worst humanitarian disaster they have witnessed.
When will our institutions start listening to these health workers and doing what is morally right?"AHW4P reiterates its calls to all medical bodies in Aotearoa and globally to not normalise the targeting of healthcare, apartheid and \textbf{genocide} with their ongoing silence.
The demands listed in @ @ @ @ @ @ @ @ @ @ immediate and permanent ceasefire in Gaza. 2 ) Call for an immediate end to the siege on Gaza. 3 ) Call for an end to the illegal occupation of Palestine by Israel. 4 ) Call for the NZ government to abide by the ICJ \textbf{ruling}. 5 ) Call for uninterrupted humanitarian aid to be allowed into Gaza. 6 ) Call for the NZ government to immediately reinstate funding to UNWRA. 7 ) Actively support the BDS campaign against the apartheid state of Israel. 8 ) Protect all healthcare workers ' right to advocate for Palestinian human rights. 9 ) Support the requests for special humanitarian visas for Palestinians in Gaza.
References 1. ' Aotearoa Healthcare Workers for Palestine open letter to all medical bodies '. https : **39 ; 919 ; TOOLONG...
Did you know \textbf{Scoop} has an Ethical Paywall?
If you ' re using \textbf{Scoop} for work, your organisation needs to pay a small \textbf{license} \textbf{fee} with \textbf{Scoop} \textbf{Pro}.
We think that ' s \textbf{fair}, because your organisation @ @ @ @ @ @ @ @ @ @, we ' ll also give your team access to \textbf{pro} news tools and keep \textbf{Scoop} free for personal use, because public access to news is important!

% matched lemmas: commit, embargo, fair, fee, genocide, license, pro, ruling, scoop, scroll
\end{textsample}
